
This is the inverse image of \(U_\ast\) under the base change
\(u:V\to U\), and hence also the inverse image of \(X_\ast\) under
the base change
\SGAThreeDiagram[0.43\textwidth]
  {figures/exp5/p271_section6_base_change_d0.png}
  {The base change defining \(V_\ast\).}
  {fig:sga3-V-6-base-change-d0}
  {Polo--Gille re-edition, Expos\'e V}

By \hyperref[par:sga3-V-3-c]{\S3.c)}, \(V_\ast\) is isomorphic to the groupoid
\(V'_\ast\), the inverse image of \(X_\ast\) under the base change
\SGAThreeDiagram[0.49\textwidth]
  {figures/exp5/p271_section6_base_change_d1.png}
  {The base change defining \(V'_\ast\).}
  {fig:sga3-V-6-base-change-d1}
  {Polo--Gille re-edition, Expos\'e V}
Thus \(V'_\ast\) admits a cokernel in \((\mathbf{Sch}/S)\). Since
\(v:V\to X_0\) is flat, surjective, and finite, it is faithfully flat
and quasi-compact, and therefore a universal effective epimorphism by
Expos\'e IV, 6.3.2.\footnote{The French re-edition mistakenly cites
Expos\'e III, which ends at \S4.}
Proposition~\hyperref[prop:sga3-V-3.1]{3.1} consequently
shows that \(X_\ast\) admits a cokernel
\(\operatorname{Coker}(d_0,d_1)\) in \((\mathbf{Sch}/S)\). We have
therefore proved the first assertion of \hyperref[lem:sga3-V-6.1-i]{part~(i)} of \hyperref[lem:sga3-V-6.1]{the following
lemma}.\footnote{\label{note:sga3-V-32}Editor's note: The argument that follows has been
slightly modified. In particular, the additional hypothesis that
\(d_0\) be flat has been moved to \hyperref[lem:sga3-V-6.1-iv]{part~(iv)}; the former part~(ii) has
been split into \hyperref[lem:sga3-V-6.1-ip]{parts~(i')} and~\hyperref[lem:sga3-V-6.1-ii]{(ii)}; and the second assertion of
\hyperref[lem:sga3-V-6.1-ip]{part~(i')} has been added.}

\medskip
\phantomsection\label{lem:sga3-V-6.1}%
\noindent\textbf{Lemma 6.1.}---\emph{Suppose that the
\((\mathbf{Sch}/S)\)-groupoid \(X_\ast\) possesses a quasi-section.
Then:}
\begin{enumerate}
\item[\SGAInlineRefTarget{lem:sga3-V-6.1:item:i}(i)]\phantomsection\label{lem:sga3-V-6.1-i}%
\emph{There exists a cokernel \((Y,p)\) of
  \((d_0,d_1)\) in \((\mathbf{Sch}/S)\). Moreover, every such
  \((Y,p)\) is a cokernel of \((d_0,d_1)\) in the category of all
  ringed spaces.}
\item[(i')]\phantomsection\label{lem:sga3-V-6.1-ip}%
\emph{The morphism \(p\) is surjective; it is open
  (respectively, universally closed) if \(d_0\) is.}
\item[(ii)]\phantomsection\label{lem:sga3-V-6.1-ii}%
\emph{Suppose that \(S\) is locally Noetherian and that
  \(X_0\) is locally of finite type (respectively, of finite type)
  over \(S\). Then \(p\) and \(Y\to S\) are locally of finite
  presentation (respectively, of finite presentation).}
\item[(iii)]\phantomsection\label{lem:sga3-V-6.1-iii}%
\emph{The morphism
  \(X_1\to X_0\times_Y X_0\) with components \(d_0\) and \(d_1\) is
  surjective.}
\item[(iv)]\phantomsection\label{lem:sga3-V-6.1-iv}%
\emph{If \((d_0,d_1)\) is an equivalence pair, then
  \(X_1\to X_0\times_Y X_0\) is an isomorphism. Moreover, if
  \(d_0:X_1\to X_0\) is flat, then \(p\) is faithfully flat.}
\end{enumerate}

Before proving the second assertion of \hyperref[lem:sga3-V-6.1-i]{part~(i)}, we shall prove
\hyperref[lem:sga3-V-6.1-ip]{parts~(i')}, \hyperref[lem:sga3-V-6.1-ii]{(ii)}, and \hyperref[lem:sga3-V-6.1-iii]{(iii)}.

\medskip
\phantomsection\label{par:sga3-V-6-a}%
\noindent\textbf{a) Proof of \hyperref[lem:sga3-V-6.1-ip]{(i')} and \hyperref[lem:sga3-V-6.1-ii]{(ii)}.}---%
We have just seen that \((Y,p)\) identifies both with
\(\operatorname{Coker}(v_0,v_1)\) and with
\(\operatorname{Coker}(u_0,u_1)\). Let \(q\) and \(r\) be the
canonical epimorphisms from \(U\) and \(V\), respectively, to \(Y\):
\SGAThreeDiagram[0.39\textwidth]
  {figures/exp5/p271_lemma61_cokernel_triangle.png}
  {The canonical epimorphisms to \(Y\).}
  {fig:sga3-V-6.1-cokernel-triangle}
  {Polo--Gille re-edition, Expos\'e V}
By hypothesis, \(v\) is surjective and finite locally free, hence
open.
